\noindent Puji dan syukur senantiasa tercurah ke hadirat Allah SWT, Tuhan Yang Maha Esa, atas segala rahmat, taufik, serta hidayah-Nya yang melimpah, sehingga penulis dapat menyelesaikan Tugas Akhir dengan judul:

\vspace{8pt}

\begin{center}
    \textbf{\large MODE KUASINORMAL DALAM RUANG-WAKTU SCHWARZSCHILD ANTI-DE SITTER}
\end{center}

\vspace{8pt}

\noindent Penulisan Tugas Akhir ini ditujukan sebagai salah satu syarat akademis dalam menyelesaikan Program Sarjana pada Departemen Fisika, Fakultas Sains dan Analitika Data, Institut Teknologi Sepuluh Nopember (FSAD-ITS). 

Penulis menyadari sepenuhnya bahwa perjalanan panjang penyusunan Tugas Akhir ini tidak luput dari tantangan dan hambatan. Namun, berkat bimbingan, dukungan, doa, serta kerja sama dari berbagai pihak, seluruh aral melintang tersebut dapat terlampaui dengan baik. Oleh karena itu, dengan segala kerendahan hati dan rasa hormat, penulis menyampaikan apresiasi serta ucapan terima kasih yang sebesar-besarnya kepada:

\begin{enumerate}
    \item Allah SWT atas segala rahmat, karunia, serta ridha-Nya sehingga penulis dapat menyelesaikan Tugas Akhir ini.

    \item Kedua orang tua dan kedua saudara, serta keluarga besar penulis atas doa, kasih sayang, dukungan moral maupun material, serta motivasi yang tiada henti selama menempuh pendidikan.

    \item Bapak Dr.rer.nat. Bintoro Anang Subagyo, M.Si., selaku dosen pembimbing, atas segala bimbingan, arahan, masukan, kesabaran, dan waktu yang telah diberikan selama proses penelitian dan penyusunan Tugas Akhir.

    \item Bapak/Ibu dosen penguji, atas kritik, saran, dan masukan yang membangun sehingga Tugas Akhir ini menjadi lebih baik.

    \item Ibu Linda Silvia, M.Si. selaku dosen wali, atas bimbingan dan arahan selama penulis menjalani studi di Departemen Fisika ITS.

    \item Seluruh dosen Departemen Fisika Fakultas Sains dan Analitika Data ITS yang telah memberikan ilmu, wawasan, dan pengalaman selama masa perkuliahan.

    \item Seluruh tenaga kependidikan Departemen Fisika Fakultas Sains dan Analitika Data ITS yang telah membantu dalam berbagai urusan administrasi selama masa studi.

    \item Teman-teman mahasiswa Fisika ITS Angkatan F40 atas kebersamaan, bantuan, semangat, serta pengalaman berharga selama menjalani perkuliahan.

    \item Rekan-rekan di Laboratorium Fisika Teoretik yang telah memberikan bantuan, diskusi, dan dukungan selama proses penelitian.

    \item Iffi Fadzilatur Rochmah yang setia menemani penulis sejak awal perkuliahan hingga penulis dapat menyelesaikan penulisan Tugas Akhir ini.

    \item Putra Naufal Tabassam, Agung Sedayu Septiawan, dan Mirza Amanda Azzahra sebagai tiga sahabat yang setia menemani penulis sejak awal hingga akhir masa perkuliahan, serta Rima Mitsaqih yang bersedia menemani dan membantu dalam penyusunan Tugas Akhir penulis.

    \item Teman-teman RUDAL Corp., atas rasa kekeluargaan, solidaritas yang erat, serta dukungan kepada penulis sepanjang berkuliah di Fisika ITS.

    \item Terakhir, terima kasih penulis tujukan kepada seluruh pihak yang telah membantu pengerjaan Tugas Akhir ini yang tidak dapat penulis sebutkan satu per satu.

\end{enumerate}



\noindent Akhir kata, penulis menyadari bahwa tulisan ini masih jauh dari kesempurnaan. Oleh karena itu, kritik dan saran yang membangun sangat diharapkan. Semoga Tugas Akhir ini dapat memberikan kontribusi nyata dan kebermanfaatan yang luas bagi perkembangan ilmu pengetahuan, khususnya di bidang fisika teoretik kosmologi.

\vspace{5em}

\begin{flushright}
    Surabaya, 26 Juni 2026
    
    \vspace{2.5cm}
    
    Ilham Rasyid Machfudi
\end{flushright}